\documentclass[11pt]{article}

\usepackage[margin=1in]{geometry}
\usepackage[T1]{fontenc}
\usepackage[utf8]{inputenc}
\usepackage{lmodern}
\usepackage{booktabs}
\usepackage{enumitem}
\usepackage{needspace}
\usepackage{parskip}
\usepackage{array}
\usepackage{tabularx}
\usepackage[hidelinks]{hyperref}

\hypersetup{
  pdftitle={bitsign MVP: Community Pilot Plan and Capability Gates},
  pdfauthor={},
  pdfsubject={Capability labels, release gates, and closed pilot protocol}
}

\newcolumntype{Y}{>{\raggedright\arraybackslash}X}
\newcolumntype{P}[1]{>{\raggedright\arraybackslash}p{#1}}
\newcommand{\code}[1]{\texttt{#1}}

\title{\textbf{bitsign MVP} \\ \large Community Pilot Plan and Capability Gates}
\author{}
\date{Draft v0.8, 22 August 2026}

\begin{document}

\maketitle

\begin{abstract}
\noindent This brief is for prospective pilot participants, Deaf ASL advisers, and
accessibility partners. For an ASL signer who prefers signing to typing, bitsign tests
whether one short signed check-in message can become useful English text for an event
staff member. The signer reviews the result and controls whether to share it; the
staff member replies through the participant's usual communication method.

The MVP combines an iPhone app with a versioned inference service. Before a closed
pilot, the build receives one of three public capability labels that state what it has
and has not demonstrated. Community review, ASL-fluent evaluation, model-provenance
review, accessibility checks, and end-to-end reliability tests are release criteria.
If the build earns a pilot-eligible label, the MVP ends with a continue, constrain,
change, or stop decision based on aggregate pilot results. Otherwise, technical review
ends with a change or stop decision without recruiting a user pilot.
\end{abstract}

\setcounter{tocdepth}{1}
\tableofcontents

\section{Scope and Governance}

This plan governs external use: what the build must demonstrate, who reviews the
evidence, when Deaf community participants may use it, and how the closed pilot runs
and terminates. Engineering and internal testing may proceed while community and
model-readiness work is underway. External recruitment begins only after every
release gate passes and the build earns a pilot-eligible capability label.

Two boundaries apply throughout:

\begin{itemize}[nosep]
  \item \textbf{Internal testing vs.\ external use.} Internal testing is gated by
        the engineering checks in R3--R5. Community readiness, model readiness, the
        capability label, and the pilot protocol additionally gate recruitment and
        external use. Internal technical readiness does not authorize recruitment.
  \item \textbf{Communication clips vs.\ research clips.} Communication clips are
        ephemeral and never train anything. Any research collection uses a separate
        informed-consent protocol with its own purpose, retention, access,
        withdrawal, and deletion rules. Pilot task flows do not collect training
        data, and pilot clips remain excluded from training and the frozen
        evaluation.
\end{itemize}

\section{Product Hypothesis}

\subsection{Primary user and use case}

The initial pilot simulates the one-way message step in an in-person community-event
check-in. A Deaf or hard-of-hearing ASL signer records one short everyday message for
an event staff member who reads English but does not understand ASL. Pilot tasks
cover arrival or registration status, schedule or location questions, and simple
non-urgent access or logistics requests. A public scope list fixes the allowed intent
classes, examples, and exclusions before evaluation. The signer can retry or discard
the output. The staff member replies through the participant's usual communication
method. (The build includes an on-device speech-to-text reply view; pilot tasks do
not use or measure it, and reverse \emph{sign generation} is outside the MVP
entirely.)

The product is an optional communication aid. Professional interpreters, qualified
translators, and established accessibility services remain the appropriate choice for
high-consequence or complex communication.

\subsection{Decisions this MVP supports}

\begin{center}
\begin{tabularx}{\textwidth}{@{}P{2.8cm}YY@{}}
\toprule
Decision & Evidence & Gate \\
\midrule
Does the model preserve meaning? & New clips from signers whose work was not used to
tune or select the model, reviewed independently by two Deaf ASL-fluent reviewers. &
The model earns a capability label in Section~\ref{sec:model}. \\
Does the complete system work? & A physical iPhone, live inference service, fixed
reliability suite, and measured warm and cold latency. & The system clears the release
gates in Section~\ref{sec:release}. \\
Is the bounded use case useful? & A closed pilot with predefined tasks, participant
feedback, observed corrections, and explicit stop conditions. & The pilot produces a
continue, constrain, change, or stop decision. \\
Can users trust the interaction? & Clear limitations, clip review, deletion behavior,
accessible status, and participant control over feedback. & The trust and data checks
pass before pilot recruitment. \\
\bottomrule
\end{tabularx}
\end{center}

A working integration establishes technical feasibility. Product validation requires
a capability that participants can understand and choose to use.

\section{Experience and Scope}

\subsection{Core flow}

\begin{center}
\code{record -> review -> translate -> read or listen -> rate or discard}
\end{center}

The signer records a 2 to 6 second ASL clip and reviews it before upload.
The service returns English text with the active capability label. Optional on-device
speech can read the result aloud. The signer can retry, discard the result, or mark it
as accurate, partly accurate, wrong, or unsafe.

For this release, \emph{live} means a short capture and response cycle. Rolling video
buffers and streaming inference require separate evaluation.

\subsection{Required product surface}

\begin{itemize}[nosep]
  \item A short introduction covering the target use case, experimental status, clip
        handling, and capability label.
  \item A framing guide that keeps the face, upper body, and both hands visible.
  \item Camera capture, clip review, reshoot, upload progress, cancel, and clear error
        states.
  \item Returned English text, optional on-device speech, retry, discard, and simple
        result feedback.
  \item A settings screen showing app and model versions plus a data-deletion contact.
  \item VoiceOver labels, Dynamic Type, sufficient contrast, and visual equivalents
        for every status communicated by sound.
\end{itemize}

Outside a consented study, categorical feedback remains on the device for the current
session and is discarded with the result. It is not linked or published. The pilot may
retain only the rating, an opaque participant code, task number, and app and model
versions for 90 days after the final session. The code key is stored separately;
withdrawal deletes linked feedback before aggregate publication, and only aggregate
feedback is published. Retaining a correction, clip, or linked result requires a
separate per-clip opt-in protocol that states its purpose, retention, access,
withdrawal, and deletion rules.

\subsection{Capability boundary}

The pilot covers ASL-to-English output from short clips only within the frozen
check-in scope list. It excludes medical, legal, financial, emergency, and other
high-consequence communication. It also excludes unrestricted continuous translation,
reverse sign generation, public App Store release, Android, and enterprise
administration. Research data collection is excluded from pilot task flows and pilot
sessions. Any later research study uses a separate informed-consent protocol and
recruits separately from the pilot.

\section{Public Commitments}

The product follows five commitments:

\begin{enumerate}[nosep]
  \item Deaf ASL signers help define the use case, evaluation rubric, pilot
        tasks, and stop conditions before the pilot begins.
  \item Public claims match the measured capability label. A successful demonstration
        clip cannot substitute for the full evaluation.
  \item Translation quality is judged primarily by ASL-fluent human review. Text
        overlap metrics remain secondary diagnostic measures.
  \item Outside a separately consented study, communication clips are used only to
        return the requested result. They never enter model training or the frozen
        model evaluation. Any research collection uses a separate opt-in protocol
        with its own stated purpose, retention, access, withdrawal, and deletion
        rules.
  \item bitsign is presented as an optional aid and never as a substitute for a
        qualified interpreter or a person's right to sign-language access.
\end{enumerate}

These commitments follow public guidance from
\href{https://gallaudet.edu/linguistics/position-statements/asl-and-ai-tools/}{Gallaudet University}
on Deaf involvement in ASL automation and from the
\href{https://wfdeaf.org/resources/consultation-on-good-practices-of-support-systems-enabling-community-inclusion-of-persons-with-disabilities/}{World Federation of the Deaf}
on community feedback and the role of professional interpreters.

\section{Model and Capability Gate}
\label{sec:model}

A capability label is a public statement, shown in the app and report, that limits
what the current build claims it can do. It is earned from the frozen evaluation, not
from a selected demonstration.

\subsection{Model provenance and deployment eligibility}

Every candidate model must have a documented source, version, training-data
provenance to the extent disclosed, and authorization for the intended pilot use.
The review is completed before external evaluation or participant recruitment. A
candidate with incomplete or incompatible deployment terms is excluded from the
pilot. Its benchmark results may inform internal research but cannot support a public
capability claim.

\subsection{Evaluation summary}

Every candidate is tested on frozen clips sampled from the public scope list and from
at least five locally recruited signers whose clips were not used to tune or select
it. Two Deaf ASL-fluent reviewers who did not contribute evaluation clips independently
compare each output with the source signing and a frozen reference. Human meaning
scores and a separate count of serious errors determine the label; text-overlap scores
are diagnostic only.

A serious error reverses or materially changes who did what, polarity, time,
quantity, or safety meaning. Results are checked both in aggregate and across signers
so a large performance gap cannot hide behind an average. Appendix~\ref{app:evaluation}
defines the frozen sets, review rubric, and reporting rules.

\subsection{Capability labels}

The thresholds are provisional. The R1 reviewers confirm them and resolve every
material concern before seeing candidate results.

\begin{center}
\begin{tabularx}{\textwidth}{@{}P{3.0cm}YY@{}}
\toprule
Label & Public claim & Pass bar \\
\midrule
Short-clip translation preview & Experimental translation of short ASL messages in
the listed community-event check-in domain. & At least 85\% of clips preserve most or
all meaning, every signer reaches at least 80\%, and no more than 2\% contain a serious
error. \\
Bounded vocabulary & Translation is limited to the phrases or isolated signs listed
in the app. & At least 90\% correct first-choice results across 125 trials, every
signer reaches at least 80\%, and no more than two trials contain a serious confusion.
\\
Technical prototype & The mobile and inference connection works, but translation
quality is unproven. & No user pilot begins. Technical review produces a change or stop
decision. \\
\bottomrule
\end{tabularx}
\end{center}

The public report gives raw aggregate counts and signer-stratified findings only when
the grouping preserves participant anonymity. The restricted evaluation record keeps
the full evidence needed to audit the decision.

\section{Technology at a Glance}

\begin{center}
\code{iOS client -> authenticated Rust API -> versioned model runtime}
\end{center}

The iOS app captures and reviews a short clip. The Rust service validates the input,
invokes the declared model, and returns text with request, model, preprocessing, and
capability identifiers. The service rejects missing models and invalid media instead
of returning placeholder output.

The first release uses direct uploads and ephemeral processing. Routine communication
clips never enter training or model-evaluation pipelines. They are deleted after
successful inference and within 24 hours after cancellation, timeout, or failure.
Study clips follow their consented protocol and remain excluded from training and the
frozen model evaluation.

Tester access is individual and revocable, transport is encrypted, and no shared
secret is embedded in the app. The service limits inputs, processing time, concurrent
work, and output length. It returns a predictable error for invalid media or overload.

The app uses SwiftUI and AVFoundation on iOS 17 or later. It stops recording after six
seconds, allows review before upload, distinguishes upload from inference progress,
and preserves the user's place through permission, network, media, and service errors.
Performance reports separate upload and model time while enforcing the end-to-end
submit-to-result gate. Appendix~\ref{app:technical} gives the detailed privacy,
security, accessibility, and performance tests.

\section{Release Gates}
\label{sec:release}

\begin{center}
\begin{tabularx}{\textwidth}{@{}P{1.3cm}P{3.4cm}Y@{}}
\toprule
Gate & Outcome & Required result \\
\midrule
R1 & Community readiness & At least three Deaf ASL signers approve the use case,
scope list, rating rubric, pilot tasks, public claims, and stop conditions. Every
material concern is resolved or the gate remains closed. \\
R2 & Model readiness & Evaluation consent, model provenance, and deployment
eligibility are recorded, the frozen evaluation is complete, and the model receives
one capability label. \\
R3 & End-to-end system & A TestFlight build on a physical iPhone sends a valid test clip to
the live service and displays the declared model's result. \\
R4 & System performance & The scripted app flow completes without a crash in 20
consecutive runs. At least 95 of 100 warm end-to-end attempts across 10 fixed test
clips succeed, and p95 from tapping Translate to rendered text is at most 12 seconds.
Invalid inputs return stable, documented errors. \\
R5 & Trust, security, and accessibility & Access, transport, resource bounds, output
bounds, and clip deletion pass their release tests. VoiceOver and Dynamic Type checks
pass, limitations are visible, and research consent is separate from communication
use. \\
\bottomrule
\end{tabularx}
\end{center}

The closed pilot requires either the short-clip translation preview or bounded
vocabulary label and all five release gates. The technical prototype label permits
technical review only and produces a change or stop decision.

\section{Closed Pilot}

Recruit 8 to 12 Deaf or hard-of-hearing ASL signers for a structured usability study.
Each participant uses the build in three predefined, low-stakes tasks that match the
capability label. At least two English-reading staff partners who do not know ASL are
assigned as evenly as the sample permits. Partners do not see the task prompt,
intended message, or reference translation. The participant's usual communication
method remains available throughout each session.

Measure:

\begin{itemize}[nosep]
  \item task completion without facilitator rescue;
  \item semantic adequacy, serious errors, retries, corrections, and abandonment;
  \item end-to-end time from recording start to a usable result;
  \item perceived usefulness compared with the participant's current method;
  \item understanding of the capability boundary and willingness to use the tool
        again in the target setting; and
  \item privacy, dignity, or trust concerns raised during the session.
\end{itemize}

A task is complete when the staff partner identifies the intended message without a
facilitator hint. A rescue occurs when the facilitator explains the model output or
redirects the interaction. Abandonment means the participant switches away before
obtaining a readable result. Serious errors use the definition in
Section~\ref{sec:model}.

Two Deaf ASL-fluent reviewers who are not pilot participants score every returned
output against the source clip and the task's predeclared intended message. The
primary serious-error count uses the first returned output for each task; retry
outputs are reported separately. The report gives raw participant, task, and output
counts alongside percentages and ranges showing statistical uncertainty. The pilot
protocol names reviewer access and the clip-deletion schedule; its clips remain
excluded from training and the frozen model evaluation.

The task-completion bar is at least 70\% without facilitator rescue. The reuse bar is
at least 75\% of participants choosing the build again in the target setting. The
safety bars require no serious error in a first returned output and no severe privacy
incident.

Any serious error pauses new sessions for adjudication. Sessions resume only after
the cause and mitigation are recorded and the affected release gates pass again.
A change to the model, prompt, preprocessing, scope list, output behavior, user flow,
or data handling starts a fresh evaluable cohort of 8 to 12 participants. An
instrumentation-only change may continue the cohort only when the R1 reviewers record
that it cannot affect the participant experience, model output, or privacy behavior.
Results from different evaluable versions are not pooled.

A severe privacy incident means unauthorized access to a clip or linked identity, or
retention beyond the documented window. New sessions pause immediately after such an
incident, a loss of deletion control, or evidence that participants misunderstand the
tool as a replacement for professional interpreting. The pilot resumes only after the
cause is corrected and the affected gates pass again; an issue that cannot be
corrected produces a stop decision.

A material pattern exists when one of the three task types completes at least 20
percentage points below the overall rate, or the same intent or capture-condition
failure affects at least three participants and the R1 reviewers confirm the pattern.
The scope cannot be narrowed by signer identity.

\Needspace{0.34\textheight}
The rules are applied from top to bottom; the first matching row is the decision:

\begin{center}
\begin{tabularx}{\textwidth}{@{}P{2.2cm}YY@{}}
\toprule
Decision & Evidence & Next step \\
\midrule
Stop & The reuse bar fails, or a safety, privacy, or scope problem cannot be corrected.
& End work on this use case and publish the reason in aggregate. \\
Change & Stop does not apply, and the task-completion, safety, or release bar fails. &
Revise the model or interaction and rerun the gates before another pilot. \\
Constrain & Both value bars and all gates pass, but a material pattern is confined to
a listed intent, task, or capture condition. & Narrow the public claim and repeat the
affected gates. \\
Continue & Both value bars and all gates pass with no material pattern. & Begin the
next bounded product iteration. \\
\bottomrule
\end{tabularx}
\end{center}

The sample supports product discovery and a next-step decision. It cannot establish
broad demand or safe use outside the target setting.

\section{Roadmap}

\begin{center}
\begin{tabularx}{\textwidth}{@{}P{3.0cm}YY@{}}
\toprule
Phase & Output & Decision \\
\midrule
1. Community definition & Target setting, scope list, user flow, review rubric, pilot
tasks, and stop conditions. & Proceed when R1 passes. \\
2. Model proof & Provenance and deployment-eligibility record, frozen evaluation,
runtime benchmark, and capability label. & Continue with a qualifying label, change
the model, or stop. \\
3. Vertical slice & Accessible TestFlight build connected to the live Rust inference
service. & Proceed when R3 through R5 pass. \\
4. Closed pilot & Aggregate task, quality, timing, usefulness, and trust results. &
Continue, constrain, change, or stop. \\
5. Public report & Capability label, evaluation method, aggregate results, known
limitations, and next-step decision. & Expand only the parts supported by evidence. \\
\bottomrule
\end{tabularx}
\end{center}

Community definition, reviewer recruitment, model-provenance review, and frozen
evaluation may run in parallel with engineering and internal testing, but they must
finish before pilot recruitment. Work beyond the pilot depends on its result. New
platforms, continuous video, reverse sign generation, and research data collection
remain outside the MVP pilot and require their own evaluation.

\section{Stop Conditions}

Recruitment or active sessions pause when documented R1 review is incomplete, a model
falls below its quality gate, a serious error awaits adjudication, or an aggregate
score hides a material gap across signers. Work also pauses when model deployment
eligibility or study consent is missing, deletion or access control fails, or a
release gate for reliability, security, accessibility, or latency no longer passes.

Work resumes only after the cause, fix, and repeated gate result are recorded. The use
case stops if usefulness fails or a safety, privacy, or scope problem cannot be
corrected.

\section{Reporting Boundary}

The public report includes the high-level architecture, evaluation method,
denominators, aggregate results, ranges showing statistical uncertainty, release-gate
status, limitations, and the next-step decision. Public build and model labels trace
to exact restricted release records.

Restricted evidence includes private service addresses, access and rate-limit
settings, capacity details, raw logs, request identifiers, evaluation manifests,
clips, reference translations, participant identities, per-signer records, detailed
operating metrics, and incident records. Public results are grouped or withheld when
a small cell could identify a participant.

\section{Public References}

This plan uses the following public sources as of 21 August 2026:

\begin{itemize}[nosep]
  \item \href{https://gallaudet.edu/linguistics/position-statements/asl-and-ai-tools/}{Gallaudet University position on ASL and AI tools};
  \item \href{https://wfdeaf.org/resources/consultation-on-good-practices-of-support-systems-enabling-community-inclusion-of-persons-with-disabilities/}{World Federation of the Deaf guidance on technology and interpreting};
  \item \href{https://aclanthology.org/2025.findings-emnlp.997/}{signer-independent evaluation in sign-language translation};
  \item \href{https://aclanthology.org/2026.lrec-1.749/}{limitations of automatic translation metrics};
  \item \href{https://developer.apple.com/accessibility/}{Apple accessibility guidance}.
\end{itemize}

\appendix
\small
\setlength{\parskip}{2pt plus 1pt minus 1pt}

\section{Evaluation Details}
\label{app:evaluation}

\subsection{Collection and separation}

Before any evaluation clip is collected, each signer receives a written protocol that
states the purpose, retention period, access roles, withdrawal and deletion process,
and publication status. Collection begins only after signed informed consent.

Each short-clip translation candidate is tested on at least 50 clips, balanced across
the intent classes in the public scope list and at least five locally recruited
signers. Their clips are created after candidate selection and withheld from all local
training, tuning, prompt selection, and model selection. For third-party weights, the
record checks disclosed provenance and says ``no known overlap'' when complete
exclusion cannot be verified.

The set and public scope list are frozen before outputs are generated and cannot be
used for later tuning. The evaluation is sentence-disjoint from local development
material and includes varied signing styles, camera distances, lighting, backgrounds,
and skin tones. This small sample does not establish demographic fairness.

\subsection{Review and reporting}

Reference translations, text-normalization rules, and rubric anchors are frozen
before candidate outputs are viewed. Two Deaf ASL-fluent reviewers who did not
contribute evaluation clips independently score each output against its source clip
and frozen reference. They do not see model identity or each other's ratings. Disputed
ratings are adjudicated before results are summarized. The report gives
pre-adjudication reviewer consistency and the number of ratings changed.

The four-point meaning rubric is:

\begin{enumerate}[nosep]
  \item intended meaning is absent or wrong;
  \item some meaning remains, but a material omission or error is present;
  \item meaning is mostly preserved without a serious error; and
  \item meaning is preserved.
\end{enumerate}

Serious errors are counted separately. Literal output, WER, BLEU, and chrF are
secondary diagnostics and cannot determine release status. The report gives raw
counts, statistical uncertainty intervals, and results by signer and capture
condition. The evaluation manifest records the exact model, preprocessing, prompt,
decoding settings, hardware, and input set.

For the bounded-vocabulary label, the supported list is frozen before testing. The
balanced set contains 25 items with one example from each of five locally recruited
signers, for 125 trials. Accuracy is macro-averaged across items and signers. The label
requires at least 90\% overall accuracy, at least 80\% for every signer, and no more
than two serious confusions.

\section{Technical Release Tests}
\label{app:technical}

\subsection{Runtime and data handling}

Short clips upload directly to the service and use ephemeral processing. Durable clip
storage, background jobs, or extra platform services are added only when measured
payload size or request duration requires them. A GPU runtime is used when the model
requires it; a CPU deployment must clear the same release gates.

Routine clips are deleted after successful inference and within 24 hours after
cancellation, timeout, or failure. Deletion tests cover on-device temporary files,
API storage, runtime caches, queues, provider logs, and crash analytics, or record that
a layer retains no clip material.

Inference telemetry uses opaque request identifiers and may record model version,
input size, timing, status, and coarse failure code. It excludes video, raw frames,
returned text, credentials, and participant identity. Separate security audit records
contain tester account, event time, authentication outcome, and source IP for up to 30
days. Access is limited to named incident responders, and audit records are linked to
inference telemetry only for a documented security investigation.

\subsection{Security and resilience}

Tester access is individual and revocable, transport is encrypted, and the app
contains no shared static secret. The service enforces fixed limits on media bytes,
duration, decoded frames, generated tokens, wall time, and concurrency. Output is
bounded plain text, and overload returns a stable error.

A versioned malformed and adversarial media fixture set runs under a fixed timeout.
Every fixture must be rejected with no process panic, hang, unexpected server error,
or residual media artifact.

\subsection{Performance and accessibility}

The scripted app flow must complete without a crash in 20 consecutive runs. The warm
performance suite runs ten end-to-end attempts for each of ten fixed test clips. A
successful attempt returns schema-valid, non-placeholder output from the declared
model within 12 seconds. At least 95 of 100 attempts must succeed, and p95 from tapping
Translate to rendered text must be at most 12 seconds.

The release record fixes the device, network profile, timeout, warm-state definition,
and percentile method, and reports both warm and cold behavior. Until a cold-start
threshold is validated, each pilot session verifies ready capacity before the first
participant task.

Accessibility checks cover VoiceOver labels and reading order, Dynamic Type, contrast,
and visual equivalents for every status communicated by sound. Consent for research
collection remains separate from ordinary communication use.

\end{document}
